\documentclass[10pt]{article}
\usepackage[T1]{fontenc}
\usepackage{lmodern,amsmath,amssymb,amsthm,xcolor}
\usepackage[margin=0.85in]{geometry}
\usepackage[colorlinks=true,allcolors=blue!60!black]{hyperref}
\newtheorem{theorem}{Theorem}
\newtheorem{corollary}[theorem]{Corollary}
\newcommand{\val}{v}
\title{Integral solutions of the Franel recurrence and sharp uniform congruences for binomial power sums}
\author{Blake Stone\\\small Developed with OpenAI Codex}
\date{November 26, 2026}
\begin{document}
\maketitle
\begin{abstract}
We prove Conjectures 4.5, 5.4 and 5.6 of Schulte, VanSchalkwyk and Yang.
Every integral solution of the classical Franel recurrence is an integer
multiple of the Franel numbers. Every other solution with integral initial
values has unbounded denominator powers at every odd prime. An entire binomial power-sum sequence avoids some prime exactly when its exponent is odd. For each
exponent $a\ge2$, we also determine the largest power of two modulo which
the sum of the $a$th powers of a binomial row is always its number of odd
entries.
\end{abstract}
\section{Initial values and denominators}
Let $F_n=\sum_{j=0}^n\binom{n}{j}^{3}$. Franel's classical recurrence,
with $F_0=1$ and $F_1=2$, reads
\begin{equation}
 n^2u_n=(7n^2-7n+2)u_{n-1}+8(n-1)^2u_{n-2},\qquad n\ge2.
 \label{eq:rec}
\end{equation}
Schulte, VanSchalkwyk and Yang \cite{report} asked in their Conjecture 4.5
which positive integer initial values yield an integer sequence. The
solution with initial values $(0,1)$ also occurs in the literature on
Franel Ap\'ery limits \cite{apery}.

\begin{theorem}
For arbitrary $u_0,u_1\in\mathbb Z$, all terms of \eqref{eq:rec} are
integers exactly when $u_1=2u_0$. Otherwise, for every odd prime $p$
and $r\ge1$, at least one of $u_{p^r-1},u_{p^r}$ has $p$-adic
valuation at most $v_p(u_1-2u_0)-2r$, with $v_p(0)=+\infty$.
\end{theorem}
\begin{proof}
Write $\delta=u_1-2u_0$. The determinant
$W_n=F_{n-1}u_n-F_nu_{n-1}$ has $W_1=\delta$.
Subtracting the two instances of \eqref{eq:rec} gives
$n^2 W_n=-8(n-1)^2W_{n-1}$, whence
\begin{equation}
 W_n=\frac{(-8)^{n-1}}{n^2}\,\delta\qquad(n\ge1).
 \label{eq:wronskian}
\end{equation}
For an odd prime and $n=p^r$, the right hand side has valuation
$v_p(\delta)-2r$ when $\delta\ne0$. The $F_n$ are integers, so two
terms with larger valuations could not give this determinant.
In particular the denominator powers are unbounded. When $\delta=0$,
uniqueness for the recurrence gives $u_n=u_0F_n$ for every $n$.
\end{proof}

\begin{corollary}
If $\delta\ne0$, then at any odd prime $p$ with $v_p(\delta)<2$,
one has $v_p(u_p)=v_p(\delta)-2$.
For $p=3$, and more generally for any odd prime such that all
$F_0,\ldots,F_{p-1}$ are units modulo $p$, one has
$v_p(u_{p^r})=v_p(\delta)-2r$ whenever $2r>v_p(\delta)$.
\end{corollary}
\begin{proof}
All terms through $u_{p-1}$ are $p$-integral by \eqref{eq:rec}. Also
$F_{p-1}\equiv\sum_{k=0}^{p-1}(-1)^{3k}=1\pmod p$.
Apply \eqref{eq:wronskian} with $n=p$.

For the second claim, Lucas's theorem applied inside the defining
binomial sum shows that $F_n\equiv\prod_i F_{d_i}\pmod p$, where
$d_i$ are the base-$p$ digits of $n$; the sum factors independently on
these digits. This identity already appears in \cite{report}.
Hence all $F_n$ are $p$-units under the hypothesis, which holds at three
since $F_0,F_1,F_2=1,2,10$. Divide \eqref{eq:wronskian} by
$F_{n-1}F_n$ and sum from $1$ to $p^r$. In the sum for
$u_{p^r}/F_{p^r}-u_0$ the unique term of least valuation is the last
one, with valuation $v_p(\delta)-2r$. It is negative by hypothesis
and cannot cancel against $u_0$.
\end{proof}

\section{Primes avoided by a whole sequence}
For $n\ge0$ and $a\ge1$ put $S_{n,a}=\sum_{k=0}^n\binom{n}{k}^{a}$.
Conjecture 5.4 of \cite{report} predicts the following classification.
\begin{theorem}
There exists a prime dividing none of $S_{0,a},S_{1,a},\ldots$
if and only if $a$ is odd.
\end{theorem}
\begin{proof}
For each odd $a$, the identity $x^a=x$ in $\mathbb F_3$ gives
$S_{n,a}\equiv\sum_k\binom nk=2^n\not\equiv0\pmod3$.
Indeed the identical argument works at any odd prime $p$ for which
$p-1\mid a-1$, by Fermat's theorem.
For even $a$ and any prime $p$, $\binom{p-1}{k}\equiv(-1)^k\pmod p$.
Thus $S_{p-1,a}\equiv\sum_{k=0}^{p-1}1=0\pmod p$, including $p=2$.
\end{proof}
The even-exponent conclusion also follows from Calkin's earlier
prime-interval divisibility theorem \cite[Proposition 3]{calkin}
with $n=p-1$ and interval parameter one. The source report gives
the modulus-three observation for the classical cubic sum.

\section{The best uniform power of two}
Let $s_2(n)$ count the ones in the binary expansion of $n$.
Conjecture 5.6 in \cite{report} asks whether $S_{n,a}/2$ is odd
exactly when $n$ is a power of two.

\begin{theorem}
The largest nonnegative integer $h$ for which
$S_{n,a}\equiv 2^{s_2(n)}\pmod{2^h}$ holds for every $n\ge1$ is
\[
 h=\begin{cases}
  2,&a\text{ odd},\\
  \min\{a,\,3+v_2(a)\},&a\text{ even}.
 \end{cases}
\]
In particular $S_{n,a}/2$ is odd exactly when $n$ is a power of two.
\end{theorem}
\begin{proof}
The coefficient of $X^k$ in $(1+X)^n$ is odd precisely when the set
bits of $k$ form a subset of those of $n$. This follows by factoring
$(1+X)^n$ modulo two using the binary expansion of $n$. Hence there
are $2^{s_2(n)}$ odd coefficients. If $n=2m$ the unpaired middle
coefficient is even since $\binom{2m}{m}=2\binom{2m-1}{m-1}$.
All odd coefficients occur in pairs of equal values $u$.
Such a pair contributes $2(u^a-1)$ to the difference from the number
of odd coefficients. All even coefficients contribute multiples of
$2^a$.

For odd $a$, each difference from an odd pair is divisible by four.
For even $a$, write $a=2^v c$ with $c$ odd. Every odd $u$ satisfies
$u^2\equiv1\pmod8$, and each further squaring raises the guaranteed
modulus by a factor of two. Thus $u^a\equiv1\pmod{2^{v+2}}$.
The extra factor two from pairing gives the asserted uniform modulus.
To see sharpness, at $n=2$ the difference is exactly $2^a$.
At $n=3$ it is $2(3^a-1)$. For odd $a$ the latter has valuation two;
for even $a$ its valuation is $v_2(a)+3$, by factoring
$3^{2^v}-1$ into $3-1$, $3+1$, and successive terms of the form
$3^{2^j}+1$ of valuation one, then observing that an odd multiplier
of the exponent preserves the valuation. Finally $h\ge2$, and
$2^{s_2(n)}$ equals two modulo four exactly when $s_2(n)=1$.
\end{proof}

\small\begin{thebibliography}{9}
\bibitem{report} A. Schulte, S. VanSchalkwyk and A. Yang,
\emph{On the divisibility and valuations of the Franel numbers},
2014 MSRI-UP research reports, pp. 127--162, Conjectures 4.5, 5.4 and 5.6.
\href{https://slmath-production-assets.sfo3.digitaloceanspaces.com/system/cms/files/81/files/original/Research_Reports_2014_MSRI-UP_%28Single_File%29.pdf}{Original report}.
\bibitem{apery} A. Straub and W. Zudilin,
\href{https://arminstraub.com/downloads/pub/franel-aperylimits.pdf}{Sums of powers
of binomials, their Ap\'ery limits, and Franel's suspicions}.
\bibitem{calkin} N. J. Calkin, \href{https://www.math.clemson.edu/~calkin/Papers/factors_sums_powers_binomial_coefficients.pdf}{Factors of sums of powers of binomial coefficients}, Acta Arith. 86 (1998), 17--26.
\end{thebibliography}
\end{document}
